%% response-to-reviewers.tex
%%
%% PLACEHOLDER point-by-point response to the IMMI reviews of
%% "Materials Science Optimization Benchmark Dataset for High-dimensional,
%%  Multi-objective, Multi-fidelity Optimization of CrabNet Hyperparameters".
%%
%% The reviewer/editor comments are transcribed verbatim; every author
%% response is an intentional [TODO] stub to be filled in after the questions
%% in author-questions.md are answered. Bracketed [Tn]/[Qn]/[En]/[Rx.y] tags
%% cross-reference review-response-plan.md.
%%
%% Self-contained: compiles with a basic TeX Live (article class + geometry,
%% xcolor, enumitem, hyperref). No sn-jnl required.
%%
%%   latexmk -pdf response-to-reviewers.tex
%%   # or: pdflatex response-to-reviewers.tex

\documentclass[11pt]{article}

\usepackage[a4paper,margin=1in]{geometry}
\usepackage[dvipsnames]{xcolor}
\usepackage{enumitem}
\usepackage{hyperref}
\hypersetup{colorlinks=true, linkcolor=RoyalBlue, urlcolor=RoyalBlue,
  citecolor=RoyalBlue}

% --- markup helpers -------------------------------------------------------- %
\definecolor{revrule}{RGB}{150,150,150}
\definecolor{respblue}{RGB}{0,80,160}

% Reviewer / editor comment: indented italic block with a grey left rule.
\newenvironment{comment}
  {\par\medskip\noindent
   \begingroup\color{revrule}\leaders\vrule width 2pt\hfill\kern0pt\par\nobreak
   \endgroup
   \begin{list}{}{\setlength{\leftmargin}{1.5em}\setlength{\rightmargin}{0.5em}}%
   \item\itshape\ignorespaces}
  {\end{list}\par}

% Author response marker.
\newcommand{\response}{\par\noindent\textbf{\color{respblue}Author response.}\ }
\newcommand{\todo}{\textbf{\color{red}[TODO]}\ }

% Cross-reference tag to the plan (e.g. \tag{T1, Q2}).
\newcommand{\tag}[1]{{\small\textcolor{respblue}{[\,#1\,]}}}

\newcommand{\reviewer}[1]{\bigskip\noindent{\large\bfseries #1}\par
  \rule{\linewidth}{0.8pt}\par}

\title{Response to Reviewers\\[-0.2em]
  \large IMMI --- CrabNet Hyperparameter Benchmark Data Descriptor}
\author{S.\,G.~Baird, J.\,N.~Parikh, T.\,D.~Sparks \\ \emph{(author list pending --- see Q11)}}
\date{\today\quad\textcolor{red}{[DRAFT --- placeholder responses]}}

\begin{document}
\maketitle

\noindent Dear Editor,\\[0.5em]
\todo Cover letter. Thank the editor and reviewers; state that the manuscript
has been revised as a Data Descriptor; summarise the headline changes in
2--4 sentences. \tag{see review-response-plan.md \S E}

\bigskip
\noindent\textbf{Summary of major changes} (placeholder --- confirm after Q1--Q14):
\begin{itemize}[nosep]
  \item \todo Scope stated explicitly as a hyperparameter-optimization benchmark. \tag{T2, Q2}
  \item \todo Sum-to-one constraint reframed and precisely defined; alloy-physics claim removed. \tag{T1, Q3}
  \item \todo Primers on CrabNet, Matbench and Ax + a researcher-workflow paragraph added. \tag{T3, Q7}
  \item \todo Data dictionary + FAIR statement added. \tag{T7, E1}
  \item \todo Objective trade-off / Pareto analysis added. \tag{T8, R1.2}
  \item \todo Figures consolidated; resolution/tick fixes; Pareto figure added. \tag{T6, Q8}
  \item \todo Multi-fidelity axes defined (and/or claim softened). \tag{T4, Q9}
  \item \todo Noise-model justification / ablation. \tag{T9, Q4}
  \item \todo Optimizer demonstration or documented downstream uses. \tag{T10, Q5}
  \item \todo Limitations paragraph + expanded related work. \tag{T11, Q10}
  \item \todo Data-accuracy corrections (repeat count; order-of-magnitude claim). \tag{T13, Q12}
  \item \todo Funding Information field completed (NSF DMR-1651668). \tag{T12, E6}
\end{itemize}

% ========================================================================= %
\reviewer{Editor}

\begin{comment}
Provide a data dictionary and demonstrate FAIR compliance: IMMI asks that data
supporting published articles meet FAIR principles and carry sufficient
metadata to describe provenance and enable reuse.
\end{comment}
\response\todo \tag{T7, E1} Point to the new data-dictionary table and FAIR
paragraph in \emph{Data Records}; give the location.

\begin{comment}
Resolve the ambiguity in the sum-to-one constraint (R1 \#1, R1 \#3, R2 \#3).
\end{comment}
\response\todo \tag{T1, E2, Q3} Summarise the reframe and the new
marginal-distribution figure; give the location.

\begin{comment}
Establish materials-community value, and make the scope explicit.
\end{comment}
\response\todo \tag{T2, E3, Q2}

\begin{comment}
Add context for the tools used: a short paragraph each on CrabNet, Matbench
(the experimental band gap task and why it was chosen), and the Ax platform;
and sketch how a typical researcher currently uses these tools, and what this
benchmark lets them do better.
\end{comment}
\response\todo \tag{T3, E4, Q7}

\begin{comment}
Figure 6 is supplied at a resolution that makes several labels and values
illegible; please supply it as vector art or at print resolution. The
logarithmic y-axis tick labelling in Figure 4 is difficult to read. Figures
1--5 are five separate single-panel histograms of similar type and could be
consolidated into one multi-panel figure.
\end{comment}
\response\todo \tag{T6, E5, Q8} Note the consolidation of Figs 1--5, the
vector Fig 6, the Fig 4 tick fix, and the added Pareto figure.

\begin{comment}
NSF DMR-1651668 is acknowledged in the text but the Funding Information field
in the submission record is blank. Please complete it.
\end{comment}
\response\todo \tag{T12, E6, Q11} Confirm the field has been completed.

% ========================================================================= %
\reviewer{Reviewer 1}

\begin{comment}
The manuscript states ``applying a contrived constraint that the sum of all
parameters must equal one'' and argues that this mimics metal-alloy composition
constraints. This requires proper justification, as the constraint is
artificial, and its physical motivation is possibly overstated.
\end{comment}
\response\todo \tag{T1, R1.1, Q3}

\begin{comment}
The dataset contains: MAE, RMSE, runtime, model size as outputs. But it is not
clear whether these exhibit meaningful trade-offs. Are there genuine Pareto
fronts, and are some objectives strongly correlated?
\end{comment}
\response\todo \tag{T8, R1.2} Report the Spearman table and Pareto fronts
(MAE\,$\leftrightarrow$\,RMSE $\rho\approx0.97$; accuracy\,$\leftrightarrow$\,runtime
$\rho\approx-0.6$; model size $\approx$ independent) and reference the new
figure.

\begin{comment}
Is Sobol sampling equivalent to conventional statistical uniform random
sampling? The authors should show whether the final 23-dimensional dataset
actually has approximately uniform marginal distributions and how the imposed
constraint affects them.
\end{comment}
\response\todo \tag{T1, R1.3, Q3} Reference the new marginal-distribution
figure (raw Sobol vs post-constraint).

\begin{comment}
To my best understanding the paper is about ML hyperparameter optimization
benchmark based on a materials dataset, rather than a materials-design
benchmark in the usual sense. This distinction should be made very clear.
\end{comment}
\response\todo \tag{T2, R1.4, Q2}

\begin{comment}
Authors should justify ``SHAP analysis as shown in Figure 6 analyzes the
difference between the true experimental output and the expected output of the
model''. SHAP values quantify the contribution of features to a model
prediction relative to a reference/expected prediction. They do not directly
measure the difference between ``true experimental output'' and expected output.
\end{comment}
\response\todo \tag{T5, R1.5, Q6} Correct the wording (attribution vs error);
confirm the fix or the figure's removal.

% ========================================================================= %
\reviewer{Reviewer 2}

\noindent\emph{Reviewer 2's assessment is a single paragraph; it is split below
into its distinct points for a structured response.}

\begin{comment}
[Novelty] The novelty and scientific depth of the contribution are limited; the
work is essentially a data generation and storage exercise, and the manuscript
reads more as a technical report describing a computational pipeline than a
research article.
\end{comment}
\response\todo \tag{T11, R2.a, Q2, Q10} Note the Data Descriptor article type
and the sharpened contribution framing.

\begin{comment}
[SHAP] The SHAP analysis in Figure 6 remains superficial and does not provide
substantial new understanding beyond what is intuitively expected (e.g., more
epochs and larger batches increase runtime).
\end{comment}
\response\todo \tag{T5, R2.b, Q6}

\begin{comment}
[Multi-fidelity] The claim of ``multi-fidelity'' is not convincingly
substantiated; the multi-fidelity dimension is not clearly defined,
demonstrated, or exploited in a manner that would distinguish this work from a
standard hyperparameter sweep.
\end{comment}
\response\todo \tag{T4, R2.c, Q9}

\begin{comment}
[Constraint] The authors themselves describe the summation-to-one constraint as
``contrived''; it does not reproduce the physics, chemistry, or thermodynamic
realities of actual compositional optimization, so the claimed bridge to
real-world materials optimization remains aspirational.
\end{comment}
\response\todo \tag{T1, R2.d, Q3}

\begin{comment}
[Noise model] The treatment of heteroskedastic noise through percentile ranks
lacks rigorous validation that it meaningfully improves surrogate fidelity
compared to simpler noise models; no quantitative comparison or ablation is
provided.
\end{comment}
\response\todo \tag{T9, R2.e, Q4}

\begin{comment}
[Optimizer demo] The practical utility of the benchmark is not convincingly
established; no demonstration is provided showing that existing adaptive design
or Bayesian optimization algorithms perform differently or more informatively
on this dataset compared to established benchmarks such as Olympus or the
particle packing benchmark previously published by the same group.
\end{comment}
\response\todo \tag{T10, R2.f, Q5}

\begin{comment}
[Scope] The scope is quite narrow, being tied exclusively to one neural network
architecture and one materials property prediction task.
\end{comment}
\response\todo \tag{T2, T11, R2.g, Q10}

\begin{comment}
[Self-assessment] The writing would benefit from more critical self-assessment
of limitations, a more thorough discussion of related work, and a clearer
articulation of what specific optimization challenges this dataset enables that
cannot already be addressed by existing resources.
\end{comment}
\response\todo \tag{T11, R2.h, Q10}

% ========================================================================= %
\reviewer{Reviewer 3}

\begin{comment}
This was a fine effort at a data-descriptor type article. A bit more context
would have been appreciated: (1) what CrabNet is, and why anyone would care
about it; (2) what Matbench is; (3) the Ax platform.
\end{comment}
\response\todo \tag{T3, R3.1--R3.3, Q7}

\begin{comment}
The readers would really benefit from getting an idea how a typical researcher
might use these tools, and how the development of this new benchmark would
improve the quality of that work (beyond what is offered in the concluding
bulleted list).
\end{comment}
\response\todo \tag{T3, T2, R3.4, Q7} Add the status-quo\,$\to$\,with-benchmark
workflow paragraph.

\end{document}
